\documentclass[10pt,a4paper]{amsart}
\usepackage[margin=19mm]{geometry}
\usepackage{amsmath,amssymb,mathtools}
\usepackage[scr=rsfs,cal=euler]{mathalfa}
\usepackage{xfrac}
\usepackage[unicode,hidelinks,bookmarks=false]{hyperref}
\newcommand{\ie}{i.e.\ }
\newcommand{\cf}{cf.\ }
\newcommand{\resp}{respectively\ }
\newcommand{\Subsection}{\subsection}
\providecommand{\nobreakdash}{\nobreak\hbox{-}}
\setlength{\emergencystretch}{2em}
\multlinegap0pt
\usepackage{bm}
\usepackage{bbm}
\usepackage{dsfont}
\usepackage{xspace}
\usepackage{cancel}
\usepackage{mathtools}
\usepackage{amsthm}
\makeatletter
\@ifundefined{thm}{\newtheorem{thm}{Theorem}}{}
\makeatother
\title[Mono-anabelian Reconstruction of Number Fields with Restrict]{Mono-anabelian Reconstruction of Number Fields with Restricted Ramification}
\author{Yu Mao, Xiao Wang}
\date{Source: arXiv:2511.05101v1 (CC BY 4.0)}
\begin{document}
\maketitle

\noindent\emph{Source and licence.} Mono-anabelian Reconstruction of Number Fields with Restricted Ramification. Version arXiv:2511.05101v1 (CC BY 4.0), distributed under the Creative Commons Attribution 4.0 International licence, as recorded in the arXiv licence metadata for this exact version and shown on the abstract page https://arxiv.org/abs/2511.05101v1. Full rights notice: licenses/arxiv-2511.05101-CC-BY-4.0.txt.\par\smallskip

\noindent\textit{Editorial scope.} Counted unit: the Thm whose source label is \texttt{None} in the article (source0.tex lines 269--280, source label \texttt{None}), delivered together with the article's own complete original proof of it (source0.tex lines 281--304, one \texttt{proof} environment of 24 lines). No mathematics is written, repaired or re-derived: statement and proof are the article's own text line for line. Required context retained uncounted: none. Every definition, display and label of those lines is retained unchanged. Omitted from this package and not redistributed: 1--268, 305--928. Nothing inside a retained excerpt is omitted or altered: every retained body is delivered exactly as the article's lines are written, so no omission receipt of any kind accompanies this package. Citations: the retained excerpts carry no citation command at all, so no bibliography entry of the article is transcribed and no reference is redistributed. Reference adaptation: none was needed and none was made; every reference inside the delivered unit resolves inside it, so no unresolved reference or citation is produced. Printed slips kept literal: no printed word, symbol, hypothesis or number of the counted statement or of its proof was altered. Layout and class: the article's own class is \texttt{article} with packages including amsmath, amsfonts, amssymb, mathrsfs, amsthm, tikz-cd, url, hyperref, todonotes; the wrapper is the lane's pinned \texttt{amsart} editorial wrapper at 10pt on A4, which restates the theorem-like environments the retained text needs and adds the article's own macro definitions that the retained text needs (none), with extra wrapper packages (bm, bbm, dsfont, xspace, cancel, mathtools). The delivered numbering is the unit's own. Assets: no retained span contains a figure environment, a graphics-inclusion command, a PDF-page inclusion, a table or a TikZ picture; no image file is copied into this package, the package declares no asset dependency and no image file is redistributed with it. Licence: version arXiv:2511.05101v1 of this article is distributed under the Creative Commons Attribution 4.0 International licence, as recorded in the arXiv licence metadata for this exact version and shown on the abstract page https://arxiv.org/abs/2511.05101v1. A full rights notice is delivered at licenses/arxiv-2511.05101-CC-BY-4.0.txt. Characters: every retained excerpt is pure ASCII, so the delivered engine, XeLaTeX with its default Latin Modern text font, reports no missing character.
\begin{thm}[Synchronisation of $\Sigma(G)$-cyclotomes]
Let $G$ be a profinite group of NF-type with density 1 restricted ramification. Let $D \in \widetilde{S}(G)$. Then the following assertions hold: \par 
(i) Let $H \subset G$ be an open subgroup, then the natural inclusion $H \hookrightarrow G$ induces the following $H$-equivariant isomorphisms
$$
\mu_{\Sigma(G)}(G) \xrightarrow{\sim} \mu_{\Sigma(H)}(H)~\text{and}~\Lambda_{\Sigma(G)}(G) \xrightarrow{\sim} \Lambda_{\Sigma(H)}(H).
$$
(ii) The natural surjection $\mathcal{I}(G,S(G)) \twoheadrightarrow k^{\times}(D)$ induces the following $D$-equivariant isomorphisms
$$
\mu_{\Sigma(G)}(G) \xrightarrow{\sim} \bigoplus_{\ell \in \Sigma(G)}\mu_{\ell^{\infty}}(D)~\text{and}~ \Lambda_{\Sigma(G)}(G) \xrightarrow{\sim} \prod_{\ell \in \Sigma(G)} \Lambda_\ell(D)
$$
where $\mu_{\ell^{\infty}}(D) := \varinjlim_n (\varinjlim_H k^{\times}(H))[\ell^n]$ for $H$ ranges over all open subgroups of $D$ and $\Lambda_\ell(D) := \varprojlim_n \mu_{\ell^{\infty}}(D)[n]$.
\end{thm}

\begin{proof}
Assertion (i) follows immediately from the fact that an open subgroup of NF-type with density 1 restricted ramification is again a profinite group of NF-type with density 1 restricted ramification, together with the fact that $\Sigma(G) = \Sigma(H)$ and Definition 3.7. \par 
Now we verify assertion (ii). Notice that to verify assertion (ii), it suffices to verify that for each $\ell \in \Sigma(G)$, there exists a $D$-equivariant isomorphism
$$
\mu_{\ell^{\infty}}(G) \xrightarrow{\sim} \mu_{\ell^{\infty}}(D)
$$
where
$$
\mu_{\ell^{\infty}}(G) := \varinjlim_n ~\mu_{\Sigma(G)}(G)[\ell^n]
$$
and
$$
\mu_{\ell^{\infty}}(D) := \varinjlim_n ~\mu(D)[\ell^n].
$$
Notice that there is a natural map
$$
\varinjlim_n ~\varinjlim_H ~\text{ker}(\mathcal{I}(H,S(H)) \to H^{\text{ab}})[\ell^n] \to \varinjlim_n ~\varinjlim_H~k^{\times}(D_H)[\ell^{n}]
$$
where $H$ ranges over all open subgroups of $G$ and $D_H := D \cap H$. In particular, we have a natural map
$$
t: \mu_{\ell^{\infty}}(G) \xrightarrow{\sim} \mu_{\ell^{\infty}}(D).
$$
The injectivity of $t$ follows from Lemma 3.4 (iii). The surjectivity of $t$ follows from the injectivity of $t$, since both $\mu_{\ell^{\infty}}(G)$ and $\mu_{\ell^{\infty}}(D)$ are abstractly isomorphic to $\mathbb{Q}_{\ell}/\mathbb{Z}_{\ell}$ and every injective endomorphism of $\mathbb{Q}_{\ell}/\mathbb{Z}_{\ell}$ is surjective. Moreover, one checks immeidately that $t$ is $D$-equivariant. This proves Theorem 3.8.
\end{proof}

\end{document}
